\documentclass[12pt,a4paper]{article}
\usepackage[utf8]{inputenc}
\usepackage[english]{babel}
\usepackage{amsmath,amsfonts,amssymb,amsthm}
\usepackage{graphicx}
\usepackage{xcolor}
\usepackage{tikz}
\usepackage{pgfplots}
\usepackage{booktabs}
\usepackage{array}
\usepackage{multirow}
\usepackage{geometry}
\usepackage{fancyhdr}
\usepackage{listings}
\usepackage{algorithm}
\usepackage{algorithmic}
\usepackage{hyperref}
\usepackage{cleveref}
\usepackage{tcolorbox}
\usepackage{enumitem}

% Code listing configuration
\lstdefinestyle{cryptoCode}{
    language=C++,
    basicstyle=\fontsize{8.5}{10}\ttfamily,
    keywordstyle=\color{blue}\bfseries,
    commentstyle=\color{green!60!black}\itshape,
    stringstyle=\color{red},
    numberstyle=\tiny\color{gray},
    numbers=left,
    numbersep=8pt,
    tabsize=2,
    showspaces=false,
    showstringspaces=false,
    showtabs=false,
    breaklines=true,
    breakatwhitespace=true,
    breakindent=20pt,
    prebreak=\raisebox{0ex}[0ex][0ex]{\ensuremath{\hookleftarrow}},
    postbreak=\mbox{\textcolor{red}{$\hookrightarrow$}\space},
    frame=tb,
    framerule=0.5pt,
    framexleftmargin=5pt,
    framexrightmargin=5pt,
    xleftmargin=10pt,
    xrightmargin=10pt,
    aboveskip=12pt,
    belowskip=12pt,
    lineskip=1pt,
    columns=flexible,
    keepspaces=true,
    captionpos=b,
    backgroundcolor=\color{gray!5}
}

\lstset{style=cryptoCode}

% Page setup
\geometry{margin=2.5cm}
\pagestyle{fancy}
\fancyhf{}
\fancyhead[L]{\textsc{Differential Cryptanalysis Guide}}
\fancyhead[R]{\textsc{CryptoGL Educational Series}}
\fancyfoot[C]{\thepage}

% Colors
\definecolor{primaryblue}{RGB}{0,102,204}
\definecolor{secondaryred}{RGB}{204,51,51}
\definecolor{accentgreen}{RGB}{51,153,51}
\definecolor{warningorange}{RGB}{255,153,51}
\definecolor{lightgray}{RGB}{245,245,245}

% Custom boxes
\newtcolorbox{definitionbox}[1]{
    colback=primaryblue!5,
    colframe=primaryblue,
    title={#1},
    fonttitle=\bfseries,
    rounded corners
}

\newtcolorbox{examplebox}[1]{
    colback=accentgreen!5,
    colframe=accentgreen,
    title={#1},
    fonttitle=\bfseries,
    rounded corners
}

\newtcolorbox{warningbox}[1]{
    colback=warningorange!5,
    colframe=warningorange,
    title={#1},
    fonttitle=\bfseries,
    rounded corners
}

\newtcolorbox{algorithmbox}[1]{
    colback=lightgray,
    colframe=black,
    title={#1},
    fonttitle=\bfseries,
    rounded corners
}

% Theorem environments
\theoremstyle{definition}
\newtheorem{definition}{Definition}[section]
\newtheorem{theorem}{Theorem}[section]
\newtheorem{lemma}{Lemma}[section]
\newtheorem{example}{Example}[section]

% Custom commands
\newcommand{\diff}{\Delta}
\newcommand{\prob}{\mathbb{P}}
\newcommand{\xor}{\oplus}
\newcommand{\Sbox}[1]{\mathsf{S}(#1)}
\newcommand{\diffchar}{\delta}
\newcommand{\round}[1]{R^{(#1)}}

\title{
    \Huge\textbf{Differential Cryptanalysis} \\
    \Large\textit{A Complete Educational Guide for Beginners} \\
    \vspace{1cm}
    % \includegraphics[width=0.3\textwidth]{differential_analysis_logo.png} % Logo can be added here if available
    \textsc{\Large CryptoGL Educational Series}
}

\author{
    \textbf{CryptoGL Educational Series} \\
    \textit{Comprehensive Cryptographic Security Analysis}
}

\date{\today}

\begin{document}

\maketitle

\begin{abstract}
This comprehensive educational guide introduces differential cryptanalysis, one of the most powerful techniques for analyzing symmetric cryptographic algorithms. Designed specifically for beginners, this document provides complete mathematical foundations, step-by-step examples, detailed algorithms, and practical defense strategies. Through clear explanations and concrete examples, readers will gain a thorough understanding of how differential cryptanalysis works, how to apply it, and how to protect against it.
\end{abstract}

\tableofcontents
\newpage

\section{Introduction to Differential Cryptanalysis}

\begin{definitionbox}{What is Differential Cryptanalysis?}
\textbf{Differential cryptanalysis} is a powerful method for analyzing symmetric-key cryptographic algorithms, particularly block ciphers. It exploits predictable differences in plaintext pairs to reveal information about the secret key by studying how these differences propagate through the cipher's rounds.
\end{definitionbox}

\subsection{Historical Context}

Differential cryptanalysis was formally introduced by Eli Biham and Adi Shamir in 1990 \cite{biham1991differential}, though the underlying principles were known earlier. Their groundbreaking work demonstrated that the Data Encryption Standard (DES) \cite{nist1999des}, previously considered secure, could be broken with $2^{47}$ operations instead of the expected $2^{56}$ brute force attack.

\subsection{Why Study Differential Cryptanalysis?}

Understanding differential cryptanalysis is crucial for several reasons:

\begin{enumerate}
    \item \textbf{Security Analysis:} It provides a systematic method to evaluate cipher strength
    \item \textbf{Cipher Design:} Knowledge helps create more secure algorithms
    \item \textbf{Attack Development:} It's a foundation for more advanced cryptanalytic techniques
    \item \textbf{Defense Strategies:} Understanding attacks enables better protection mechanisms
\end{enumerate}

\section{Mathematical Foundations}

\subsection{Basic Concepts and Notation}

Before diving into differential cryptanalysis, we need to establish the mathematical framework.

\begin{definitionbox}{XOR Operation and Differences}
The XOR operation ($\oplus$) is fundamental to differential cryptanalysis. For two bit strings $a$ and $b$ of equal length:
\begin{align}
a \oplus b = c
\end{align}
where $c_i = a_i \oplus b_i$ for each bit position $i$.

The \textbf{difference} between two values $x$ and $x'$ is defined as:
\begin{align}
\diff x = x \oplus x'
\end{align}
\end{definitionbox}

\begin{examplebox}{XOR Difference Example}
Let's compute the difference between two 8-bit values:
\begin{align}
x &= 10110011_2 = \text{0xB3} \\
x' &= 11010001_2 = \text{0xD1} \\
\diff x &= x \oplus x' = 10110011 \oplus 11010001 = 01100010_2 = \text{0x62}
\end{align}
\end{examplebox}

\subsection{Properties of XOR Differences}

The XOR operation has several important properties that make it ideal for differential analysis:

\begin{theorem}[XOR Difference Properties]
For any values $a$, $b$, and $c$:
\begin{enumerate}
    \item \textbf{Self-inverse:} $a \oplus a = 0$
    \item \textbf{Commutative:} $a \oplus b = b \oplus a$
    \item \textbf{Associative:} $(a \oplus b) \oplus c = a \oplus (b \oplus c)$
    \item \textbf{Linear:} $(a \oplus b) \oplus (c \oplus d) = (a \oplus c) \oplus (b \oplus d)$
\end{enumerate}
\end{theorem}

\begin{examplebox}{Linearity Example}
The linearity property is crucial for differential analysis. If we have:
\begin{align}
\diff x &= x_1 \oplus x_2 \\
\diff y &= y_1 \oplus y_2
\end{align}
Then for any linear function $f$:
\begin{align}
f(\diff x) = f(x_1 \oplus x_2) = f(x_1) \oplus f(x_2) = \diff f(x)
\end{align}
\end{examplebox}

\subsection{Differential Probability}

\begin{definitionbox}{Differential Probability}
For a function $f: \{0,1\}^n \rightarrow \{0,1\}^m$ and input/output differences $\alpha$ and $\beta$, the \textbf{differential probability} is:
\begin{align}
\prob(\alpha \rightarrow \beta) = \frac{|\{x : f(x) \oplus f(x \oplus \alpha) = \beta\}|}{2^n}
\end{align}
This represents the probability that input difference $\alpha$ produces output difference $\beta$ through function $f$.
\end{definitionbox}

\section{Core Concepts of Differential Cryptanalysis}

\subsection{Differential Characteristics}

\begin{definitionbox}{Differential Characteristic}
A \textbf{differential characteristic} is a sequence of differences through multiple rounds of a cipher:
\begin{align}
\diffchar: \diff P \rightarrow \diff C_1 \rightarrow \diff C_2 \rightarrow \cdots \rightarrow \diff C_r
\end{align}
where $\diff P$ is the plaintext difference, $\diff C_i$ is the difference after round $i$, and $r$ is the number of rounds.

The \textbf{probability} of a characteristic is:
\begin{align}
\prob(\diffchar) = \prod_{i=1}^{r} \prob(\diff C_{i-1} \rightarrow \diff C_i)
\end{align}
\end{definitionbox}

A differential characteristic represents a specific path that differences follow through the cipher's rounds. Each step in this path has an associated probability, and the overall probability is the product of all individual step probabilities.

\subsubsection{Types of Differential Characteristics}

\begin{enumerate}
    \item \textbf{Deterministic Characteristics:} Have probability 1.0 (very rare)
    \item \textbf{High-Probability Characteristics:} Probability significantly above random
    \item \textbf{Iterative Characteristics:} Repeat the same pattern over multiple rounds
    \item \textbf{Truncated Characteristics:} Specify only partial difference patterns
\end{enumerate}

\begin{examplebox}{3-Round Characteristic Example}
Consider a simple 3-round characteristic for a 16-bit cipher:

\textbf{Round 0 (Plaintext):} $\diff P = \text{0x1000}$
\begin{align}
\text{Probability: } 1.0 \text{ (given input difference)}
\end{align}

\textbf{Round 1:} $\diff C_1 = \text{0xA000}$
\begin{align}
\prob(\text{0x1000} \rightarrow \text{0xA000}) = \frac{2}{16} = 0.125
\end{align}

\textbf{Round 2:} $\diff C_2 = \text{0x6000}$
\begin{align}
\prob(\text{0xA000} \rightarrow \text{0x6000}) = \frac{2}{16} = 0.125
\end{align}

\textbf{Round 3:} $\diff C_3 = \text{0x4000}$
\begin{align}
\prob(\text{0x6000} \rightarrow \text{0x4000}) = \frac{4}{16} = 0.25
\end{align}

\textbf{Overall Characteristic Probability:}
\begin{align}
\prob(\diffchar) = 1.0 \times 0.125 \times 0.125 \times 0.25 = \frac{1}{256}
\end{align}

This means approximately 1 in every 256 plaintext pairs with difference 0x1000 will follow this exact characteristic path.
\end{examplebox}

\subsubsection{Characteristic Quality Metrics}

\begin{enumerate}
    \item \textbf{Probability:} Higher probability means fewer pairs needed
    \item \textbf{Signal-to-Noise Ratio:} How well the characteristic stands out from random
    \item \textbf{Key Recovery Potential:} How much key information can be recovered
    \item \textbf{Data Complexity:} Number of plaintext pairs required
\end{enumerate}

\begin{examplebox}{Data Complexity Calculation}
For a characteristic with probability $p$, the expected number of pairs needed is:
\begin{align}
N_{\text{pairs}} = \frac{c}{p}
\end{align}
where $c$ is a constant depending on the desired success probability (typically $c \approx 3$ to $8$).

For our example characteristic with $p = \frac{1}{256}$:
\begin{align}
N_{\text{pairs}} = \frac{4}{\frac{1}{256}} = 4 \times 256 = 1024 \text{ pairs}
\end{align}
\end{examplebox}

\subsubsection{Finding Good Characteristics}

The process of finding effective differential characteristics involves:

\begin{enumerate}
    \item \textbf{Exhaustive Search:} Try all possible difference patterns (limited scope)
    \item \textbf{Heuristic Search:} Use algorithms like branch-and-bound
    \item \textbf{Automated Tools:} SAT solvers, MILP optimization
    \item \textbf{Pattern Recognition:} Identify recurring structures
\end{enumerate}

\subsection{The Substitution Box (S-box) Analysis}

S-boxes are the primary source of nonlinearity in most block ciphers, making them crucial for differential analysis.

\begin{examplebox}{Simple 4-bit S-box Analysis}
Consider a 4-bit S-box with the following mapping:
\begin{center}
\begin{tabular}{c|cccccccccccccccc}
Input  & 0 & 1 & 2 & 3 & 4 & 5 & 6 & 7 & 8 & 9 & A & B & C & D & E & F \\
\hline
Output & E & 4 & D & 1 & 2 & F & B & 8 & 3 & A & 6 & C & 5 & 9 & 0 & 7 \\
\end{tabular}
\end{center}

Let's analyze the differential $\diff x = 1 \rightarrow \diff y$:
\begin{align}
\Sbox{0} \oplus \Sbox{1} &= \text{0xE} \oplus \text{0x4} = \text{0xA} \\
\Sbox{2} \oplus \Sbox{3} &= \text{0xD} \oplus \text{0x1} = \text{0xC} \\
\Sbox{4} \oplus \Sbox{5} &= \text{0x2} \oplus \text{0xF} = \text{0xD} \\
&\vdots
\end{align}

The differential distribution table (DDT) shows how often each input difference produces each output difference.
\end{examplebox}

\subsection{Building the Differential Distribution Table}

The Differential Distribution Table (DDT) is a fundamental tool for analyzing the differential properties of S-boxes. It systematically records how input differences map to output differences across all possible inputs.

\begin{definitionbox}{Differential Distribution Table (DDT)}
For an S-box $S: \{0,1\}^n \rightarrow \{0,1\}^m$, the DDT is a $2^n \times 2^m$ table where:
\begin{align}
DDT[\diff x][\diff y] = |\{x : S(x) \oplus S(x \oplus \diff x) = \diff y\}|
\end{align}

This counts how many input values $x$ produce output difference $\diff y$ when the input difference is $\diff x$.
\end{definitionbox}

\subsubsection{Step-by-Step DDT Construction}

\begin{algorithmbox}{Algorithm 1: Computing Differential Distribution Table}
\begin{algorithmic}[1]
\REQUIRE S-box $S$ with $n$-bit input and $m$-bit output
\ENSURE Differential Distribution Table $DDT$
\STATE Initialize $DDT[2^n][2^m]$ to all zeros
\FOR{$\diff x = 0$ to $2^n - 1$}
    \FOR{$x = 0$ to $2^n - 1$}
        \STATE $x' = x \oplus \diff x$
        \STATE $\diff y = S[x] \oplus S[x']$
        \STATE $DDT[\diff x][\diff y] = DDT[\diff x][\diff y] + 1$
    \ENDFOR
\ENDFOR
\RETURN $DDT$
\end{algorithmic}
\end{algorithmbox}

\begin{examplebox}{Manual DDT Construction Example}
Let's build the DDT for a simple 3-bit S-box:
\begin{center}
\begin{tabular}{c|cccccccc}
$x$ & 0 & 1 & 2 & 3 & 4 & 5 & 6 & 7 \\
\hline
$S(x)$ & 6 & 4 & 2 & 5 & 7 & 1 & 3 & 0 \\
\end{tabular}
\end{center}

\textbf{Step 1:} Calculate for input difference $\diff x = 1$
\begin{align}
x = 0: \quad & S(0) \oplus S(0 \oplus 1) = S(0) \oplus S(1) = 6 \oplus 4 = 2 \\
x = 1: \quad & S(1) \oplus S(1 \oplus 1) = S(1) \oplus S(0) = 4 \oplus 6 = 2 \\
x = 2: \quad & S(2) \oplus S(2 \oplus 1) = S(2) \oplus S(3) = 2 \oplus 5 = 7 \\
x = 3: \quad & S(3) \oplus S(3 \oplus 1) = S(3) \oplus S(2) = 5 \oplus 2 = 7 \\
x = 4: \quad & S(4) \oplus S(4 \oplus 1) = S(4) \oplus S(5) = 7 \oplus 1 = 6 \\
x = 5: \quad & S(5) \oplus S(5 \oplus 1) = S(5) \oplus S(4) = 1 \oplus 7 = 6 \\
x = 6: \quad & S(6) \oplus S(6 \oplus 1) = S(6) \oplus S(7) = 3 \oplus 0 = 3 \\
x = 7: \quad & S(7) \oplus S(7 \oplus 1) = S(7) \oplus S(6) = 0 \oplus 3 = 3
\end{align}

\textbf{Results for $\diff x = 1$:}
\begin{itemize}
    \item Output difference 2 appears 2 times: $DDT[1][2] = 2$
    \item Output difference 7 appears 2 times: $DDT[1][7] = 2$
    \item Output difference 6 appears 2 times: $DDT[1][6] = 2$
    \item Output difference 3 appears 2 times: $DDT[1][3] = 2$
\end{itemize}
\end{examplebox}

\subsubsection{Complete DDT Analysis}

\begin{examplebox}{Complete 3-bit S-box DDT}
After computing all input differences, we get the complete DDT:

\begin{center}
\begin{tabular}{c|cccccccc}
$\diff x \backslash \diff y$ & 0 & 1 & 2 & 3 & 4 & 5 & 6 & 7 \\
\hline
0 & 8 & 0 & 0 & 0 & 0 & 0 & 0 & 0 \\
1 & 0 & 0 & 2 & 2 & 0 & 0 & 2 & 2 \\
2 & 0 & 2 & 0 & 2 & 2 & 2 & 0 & 0 \\
3 & 0 & 2 & 2 & 0 & 2 & 2 & 0 & 0 \\
4 & 0 & 0 & 2 & 2 & 0 & 0 & 2 & 2 \\
5 & 0 & 2 & 0 & 2 & 2 & 2 & 0 & 0 \\
6 & 0 & 2 & 2 & 0 & 2 & 2 & 0 & 0 \\
7 & 0 & 0 & 2 & 2 & 0 & 0 & 2 & 2 \\
\end{tabular}
\end{center}

\textbf{Key Observations:}
\begin{enumerate}
    \item Row 0 has entry 8 in column 0: input difference 0 always produces output difference 0
    \item Maximum entry (excluding row 0) is 2, giving probability $\frac{2}{8} = \frac{1}{4}$
    \item Each row sums to 8 (total number of possible inputs)
    \item No entry is greater than 4 (good differential property)
\end{enumerate}
\end{examplebox}

\subsubsection{DDT Properties and Analysis}

\begin{theorem}[DDT Properties]
For any DDT of an $n$-bit to $m$-bit S-box:
\begin{enumerate}
    \item \textbf{Row Sum:} $\sum_{j=0}^{2^m-1} DDT[i][j] = 2^n$ for all $i$
    \item \textbf{First Row:} $DDT[0][j] = 0$ for $j \neq 0$, and $DDT[0][0] = 2^n$
    \item \textbf{Symmetry:} If $S$ is a permutation, then column sums equal row sums
    \item \textbf{Maximum Entry:} For good cryptographic properties, $\max_{i \neq 0, j} DDT[i][j] \leq 2^{n-2}$
\end{enumerate}
\end{theorem}

\begin{examplebox}{Differential Probability Calculation}
From the DDT, we can directly calculate differential probabilities:

\textbf{Example 1:} $\prob(1 \rightarrow 2) = \frac{DDT[1][2]}{2^n} = \frac{2}{8} = \frac{1}{4}$

\textbf{Example 2:} $\prob(3 \rightarrow 5) = \frac{DDT[3][5]}{2^n} = \frac{2}{8} = \frac{1}{4}$

\textbf{Example 3:} $\prob(2 \rightarrow 0) = \frac{DDT[2][0]}{2^n} = \frac{0}{8} = 0$ (impossible)

The maximum differential probability for this S-box is $\frac{1}{4}$, which is considered acceptable for a 3-bit S-box.
\end{examplebox}

\subsubsection{Cryptographic Significance}

\begin{warningbox}{Security Implications}
\textbf{High DDT Entries:} Large values in the DDT indicate predictable differential behavior, making the S-box vulnerable to differential attacks.

\textbf{Security Threshold:} For an $n$-bit S-box, entries should not exceed $2^{n-2}$:
\begin{itemize}
    \item 4-bit S-box: maximum entry $\leq$ 4
    \item 6-bit S-box: maximum entry $\leq$ 16  
    \item 8-bit S-box: maximum entry $\leq$ 64
\end{itemize}

\textbf{Random S-box:} A randomly chosen S-box typically has much higher maximum entries, making designed S-boxes preferable for cryptographic use.
\end{warningbox}

\subsubsection{Implementation Considerations}

\begin{algorithmbox}{Optimized DDT Construction}
\begin{algorithmic}[1]
\REQUIRE S-box $S[0..2^n-1]$
\ENSURE DDT with maximum differential probability
\STATE Initialize $DDT[2^n][2^n]$ to zeros
\STATE $max\_prob = 0$
\FOR{$dx = 1$ to $2^n - 1$} \COMMENT{Skip $dx = 0$}
    \FOR{$x = 0$ to $2^n - 1$}
        \STATE $dy = S[x] \oplus S[x \oplus dx]$
        \STATE $DDT[dx][dy] = DDT[dx][dy] + 1$
        \STATE $max\_prob = \max(max\_prob, DDT[dx][dy])$
    \ENDFOR
\ENDFOR
\STATE $max\_differential\_probability = \frac{max\_prob}{2^n}$
\RETURN $DDT, max\_differential\_probability$
\end{algorithmic}
\end{algorithmbox}

\subsubsection{Using DDT for Cryptanalysis}

The DDT enables systematic analysis of cipher components:

\begin{enumerate}
    \item \textbf{Characteristic Construction:} Find high-probability paths through multiple rounds
    \item \textbf{Attack Complexity Estimation:} Calculate required plaintext pairs
    \item \textbf{Security Assessment:} Evaluate resistance to differential attacks
    \item \textbf{S-box Comparison:} Compare different S-box designs objectively
\end{enumerate}

\begin{examplebox}{Multi-Round Analysis}
For a 2-round cipher using our example S-box:

\textbf{Round 1:} Best differential $1 \rightarrow 2$ with probability $\frac{1}{4}$
\textbf{Round 2:} Continue with $2 \rightarrow 7$ with probability $\frac{1}{4}$

\textbf{Combined:} $1 \rightarrow 2 \rightarrow 7$ with probability $\frac{1}{4} \times \frac{1}{4} = \frac{1}{16}$

\textbf{Attack Feasibility:} Need approximately $16 \times 4 = 64$ plaintext pairs for reliable detection.
\end{examplebox}

\begin{examplebox}{Complete DDT Example}
For our 4-bit S-box example, here's a portion of the DDT:

\begin{center}
\begin{tabular}{c|cccccccc}
$\diff x \backslash \diff y$ & 0 & 1 & 2 & 3 & 4 & 5 & 6 & 7 \\
\hline
0 & 16 & 0 & 0 & 0 & 0 & 0 & 0 & 0 \\
1 & 0 & 0 & 2 & 0 & 2 & 2 & 0 & 2 \\
2 & 0 & 2 & 0 & 2 & 0 & 0 & 4 & 0 \\
3 & 0 & 0 & 0 & 2 & 2 & 0 & 2 & 2 \\
$\vdots$ & $\vdots$ & $\vdots$ & $\vdots$ & $\vdots$ & $\vdots$ & $\vdots$ & $\vdots$ & $\vdots$ \\
\end{tabular}
\end{center}

The entry $DDT[1][2] = 2$ means that input difference 1 produces output difference 2 for exactly 2 out of 16 possible inputs, giving probability $\frac{2}{16} = \frac{1}{8}$.
\end{examplebox}

\section{Step-by-Step Differential Attack Methodology}

\subsection{Phase 1: Characteristic Discovery}

The first phase involves finding good differential characteristics through the cipher.

\begin{algorithmbox}{Algorithm 2: Finding Differential Characteristics}
\begin{algorithmic}[1]
\REQUIRE Cipher structure, S-boxes, number of rounds $r$
\ENSURE Set of differential characteristics with high probability
\STATE Choose initial plaintext difference $\diff P$
\STATE Initialize characteristic list $\mathcal{C} = \emptyset$
\FOR{each possible round 1 output difference $\diff C_1$}
    \IF{$\prob(\diff P \rightarrow \diff C_1) > \text{threshold}$}
        \STATE Recursively extend characteristic to round 2, 3, ..., $r-1$
        \STATE Add completed characteristics to $\mathcal{C}$
    \ENDIF
\ENDFOR
\STATE Sort $\mathcal{C}$ by probability (highest first)
\RETURN Top characteristics from $\mathcal{C}$
\end{algorithmic}
\end{algorithmbox}

\subsection{Phase 2: Plaintext Pair Generation}

\begin{examplebox}{Plaintext Pair Generation Example}
Suppose we have a differential characteristic:
\begin{align}
\diff P = \text{0x40800000} \rightarrow \diff C_{r-1} = \text{0x04000000}
\end{align}

We generate plaintext pairs $(P_1, P_2)$ such that $P_1 \oplus P_2 = \text{0x40800000}$:

\begin{center}
\begin{tabular}{ccc}
\textbf{Pair} & $P_1$ & $P_2 = P_1 \oplus \text{0x40800000}$ \\
\hline
1 & 0x12345678 & 0x52B45678 \\
2 & 0xABCDEF00 & 0xEB4DEF00 \\
3 & 0x11111111 & 0x51911111 \\
$\vdots$ & $\vdots$ & $\vdots$ \\
\end{tabular}
\end{center}
\end{examplebox}

\subsection{Phase 3: Ciphertext Analysis}

\begin{algorithmbox}{Algorithm 3: Key Recovery from Differential Pairs}
\begin{algorithmic}[1]
\REQUIRE Plaintext pairs $(P_i, P_i')$, corresponding ciphertext pairs $(C_i, C_i')$
\REQUIRE Differential characteristic $\diffchar$ with probability $p$
\ENSURE Candidate subkeys for the last round
\STATE Initialize key counter array $\text{Count}[2^{k}]$ to zeros, where $k$ is subkey length
\FOR{each plaintext-ciphertext pair $(P_i, P_i', C_i, C_i')$}
    \FOR{each possible last round subkey $K_r$}
        \STATE Partially decrypt: $X_i = \text{InverseRound}(C_i, K_r)$
        \STATE Partially decrypt: $X_i' = \text{InverseRound}(C_i', K_r)$
        \STATE Compute difference: $\diff X = X_i \oplus X_i'$
        \IF{$\diff X$ matches expected difference from characteristic}
            \STATE $\text{Count}[K_r] = \text{Count}[K_r] + 1$
        \ENDIF
    \ENDFOR
\ENDFOR
\STATE Find subkeys with count significantly above expected random count
\RETURN Top candidate subkeys
\end{algorithmic}
\end{algorithmbox}

\section{Complete Worked Example: Mini-Cipher Attack}

Let's work through a complete differential attack on a simplified cipher to illustrate all concepts.

\subsection{Mini-Cipher Specification}

Our educational cipher has the following properties:
\begin{itemize}
    \item Block size: 16 bits
    \item Key size: 16 bits  
    \item Number of rounds: 3
    \item S-box: 4-bit to 4-bit (same as previous example)
    \item Round function: SubBytes $\rightarrow$ ShiftRows $\rightarrow$ AddRoundKey
\end{itemize}

\begin{examplebox}{Mini-Cipher Round Function}
For a 16-bit block split into four 4-bit nibbles $[n_0, n_1, n_2, n_3]$:

\textbf{SubBytes:} Apply S-box to each nibble
\begin{align}
[n_0, n_1, n_2, n_3] \rightarrow [S(n_0), S(n_1), S(n_2), S(n_3)]
\end{align}

\textbf{ShiftRows:} Swap positions of nibbles 1 and 3
\begin{align}
[s_0, s_1, s_2, s_3] \rightarrow [s_0, s_3, s_2, s_1]
\end{align}

\textbf{AddRoundKey:} XOR with round key
\begin{align}
[r_0, r_1, r_2, r_3] \rightarrow [r_0 \oplus k_0, r_1 \oplus k_1, r_2 \oplus k_2, r_3 \oplus k_3]
\end{align}
\end{examplebox}

\subsection{Step 1: Finding a Good Differential Characteristic}

We need to find a characteristic that works through 2 rounds (attacking the 3rd round).

\begin{examplebox}{Characteristic Discovery Process}
Let's try input difference $\diff P = \text{0x1000}$ (difference only in the first nibble):

\textbf{Round 1:}
\begin{align}
\diff P &= \text{0x1000} = [1, 0, 0, 0] \\
\text{After SubBytes:} &\quad \text{Look up } DDT[1][\diff y] \text{ for all } \diff y \\
\text{High probability:} &\quad [1, 0, 0, 0] \rightarrow [A, 0, 0, 0] \text{ with prob } \frac{2}{16} \\
\text{After ShiftRows:} &\quad [A, 0, 0, 0] \rightarrow [A, 0, 0, 0] \\
\text{After AddRoundKey:} &\quad [A, 0, 0, 0] \text{ (XOR doesn't change differences)}
\end{align}

\textbf{Round 2:}
\begin{align}
\text{Input:} &\quad [A, 0, 0, 0] \\
\text{After SubBytes:} &\quad [DDT[A][\text{best}], 0, 0, 0] \\
\text{Continue analysis...}
\end{align}

After complete analysis, we find:
\begin{align}
\text{0x1000} \rightarrow \text{0xA000} \rightarrow \text{0x6000}
\end{align}
with probability $\frac{2}{16} \times \frac{2}{16} = \frac{1}{64}$.
\end{examplebox}

\subsection{Step 2: Collecting Plaintext-Ciphertext Pairs}

\begin{examplebox}{Data Collection}
We need approximately $\frac{64}{\text{signal-to-noise ratio}}$ pairs. For reliable detection, let's collect 1000 pairs.

Generate pairs $(P_i, P_i')$ where $P_i \oplus P_i' = \text{0x1000}$:

\begin{center}
\begin{tabular}{cccc}
\textbf{Pair} & $P_i$ & $P_i'$ & $\diff P$ \\
\hline
1 & 0x2847 & 0x3847 & 0x1000 \\
2 & 0x9A3C & 0x8A3C & 0x1000 \\
3 & 0x7F91 & 0x6F91 & 0x1000 \\
$\vdots$ & $\vdots$ & $\vdots$ & $\vdots$ \\
1000 & 0x1234 & 0x0234 & 0x1000 \\
\end{tabular}
\end{center}

Encrypt each pair to get $(C_i, C_i')$.
\end{examplebox}

\subsection{Step 3: Key Recovery Attack}

\begin{examplebox}{Last Round Key Recovery}
For each ciphertext pair $(C_i, C_i')$ and each possible last round key $K_3$:

\begin{enumerate}
    \item Partially decrypt: $X_i = C_i \oplus K_3$, $X_i' = C_i' \oplus K_3$
    \item Apply inverse ShiftRows: $Y_i = \text{InvShiftRows}(X_i)$, $Y_i' = \text{InvShiftRows}(X_i')$  
    \item Check if $Y_i \oplus Y_i'$ has the expected pattern from our characteristic
\end{enumerate}

\textbf{Key Counting:}
\begin{center}
\begin{tabular}{cc}
\textbf{Key Candidate} & \textbf{Count} \\
\hline
0x2A4F & 387 \\
0x1B3E & 3 \\
0x7C91 & 2 \\
0x4D82 & 1 \\
$\vdots$ & $\vdots$ \\
\end{tabular}
\end{center}

The key 0x2A4F has significantly more hits than expected by chance ($\approx$15), indicating it's likely the correct key!
\end{examplebox}

\section{Advanced Differential Techniques}

\subsection{Truncated Differentials}

\begin{definitionbox}{Truncated Differentials}
A \textbf{truncated differential} \cite{knudsen1994truncated} specifies differences in only part of the state, allowing more flexibility in the attack. Instead of specifying exact bit differences, we specify patterns like:
\begin{align}
(\text{any}, 0, \text{any}, 0) \rightarrow (0, \text{any}, 0, \text{any})
\end{align}
where "any" means any non-zero difference and "0" means zero difference.
\end{definitionbox}

\subsection{Higher-Order Differentials}

\begin{definitionbox}{Higher-Order Differentials}
For a function $f$, the \textbf{$k$-th order differential} is defined as:
\begin{align}
\Delta^k_{\alpha_1,\ldots,\alpha_k} f(x) = \bigoplus_{\substack{S \subseteq \{\alpha_1,\ldots,\alpha_k\}}} f\left(x \oplus \bigoplus_{\alpha_i \in S} \alpha_i\right)
\end{align}

This generalizes first-order differentials and can attack ciphers resistant to classical differential analysis.
\end{definitionbox}

\subsection{Impossible Differentials}

\begin{definitionbox}{Impossible Differentials}
An \textbf{impossible differential} \cite{biham1999impossible} is a differential characteristic with probability exactly zero. If we observe a plaintext-ciphertext pair following an impossible differential for a particular key guess, we can eliminate that key as impossible.
\end{definitionbox}

\section{Defense Mechanisms Against Differential Cryptanalysis}

\subsection{S-box Design Criteria}

To resist differential attacks, S-boxes should satisfy several criteria:

\begin{enumerate}
    \item \textbf{Low Maximum Differential Probability:}
    \begin{align}
    \max_{\alpha \neq 0, \beta} \prob(\alpha \rightarrow \beta) \leq \frac{1}{2^{n-2}}
    \end{align}
    where $n$ is the S-box input size.
    
    \item \textbf{No Fixed Points:} $S(x) \neq x$ for all $x$
    
    \item \textbf{No Opposite Fixed Points:} $S(x) \neq \overline{x}$ for all $x$
    
    \item \textbf{Bijective:} Each output appears exactly once (for invertibility)
\end{enumerate}

\begin{examplebox}{AES S-box Resistance}
The AES S-box achieves maximum differential probability of $\frac{4}{256} = \frac{1}{64}$, which is close to optimal for an 8-bit S-box. This makes differential attacks on full AES computationally infeasible.
\end{examplebox}

\subsection{Linear Layer Design}

The linear layer (like MixColumns in AES) should provide:

\begin{enumerate}
    \item \textbf{Maximum Distance Separable (MDS):} Ensures optimal diffusion
    \item \textbf{Branch Number:} High branch number increases the minimum number of active S-boxes
    \item \textbf{Wide Trail Strategy:} Makes it difficult to find low-weight characteristics
\end{enumerate}

\begin{definitionbox}{Branch Number}
For a linear transformation $L$, the \textbf{branch number} is:
\begin{align}
\mathcal{B}(L) = \min_{x \neq 0} \{\text{wt}(x) + \text{wt}(L(x))\}
\end{align}
where $\text{wt}(\cdot)$ is the Hamming weight (number of non-zero elements).
\end{definitionbox}

\subsection{Round Number Selection}

The number of rounds must be sufficient to:

\begin{enumerate}
    \item Make the best differential characteristics have negligible probability
    \item Ensure multiple applications of the round function eliminate any structure
    \item Provide adequate security margin beyond the theoretical minimum
\end{enumerate}

\begin{algorithmbox}{Algorithm 4: Security Margin Calculation}
\begin{algorithmic}[1]
\REQUIRE Cipher specifications, S-box properties
\ENSURE Minimum number of rounds for security
\STATE Find best differential characteristic probability $p_{\text{best}}$
\STATE Calculate theoretical attack complexity: $C_{\text{theory}} = \frac{1}{p_{\text{best}}}$
\STATE Determine practical attack complexity considering:
\STATE \quad - Data complexity
\STATE \quad - Time complexity  
\STATE \quad - Memory requirements
\STATE Set minimum rounds such that $C_{\text{practical}} > 2^{\text{key size}}$
\STATE Add security margin: $r_{\text{secure}} = r_{\text{minimum}} + \text{margin}$
\RETURN $r_{\text{secure}}$
\end{algorithmic}
\end{algorithmbox}

\section{Practical Implementation Considerations}

\subsection{Data Collection Strategies}

\begin{warningbox}{Adaptive vs. Non-Adaptive Attacks}
\textbf{Non-adaptive attacks:} Choose all plaintext pairs before seeing any ciphertext. More realistic in practice.

\textbf{Adaptive attacks:} Choose subsequent plaintexts based on previous ciphertext observations. More powerful but often impractical.

Most real-world scenarios require non-adaptive approaches due to limited oracle access.
\end{warningbox}

\subsection{Statistical Analysis}

\begin{examplebox}{Chi-Square Test for Key Detection}
To determine if a key candidate is correct, use the chi-square test:

\begin{align}
\chi^2 = \sum_{i} \frac{(\text{observed}_i - \text{expected}_i)^2}{\text{expected}_i}
\end{align}

For a key counting array with expected uniform distribution:
\begin{align}
\text{expected count} &= \frac{\text{total pairs}}{\text{number of possible keys}} \\
\chi^2 &= \sum_{k} \frac{(\text{Count}[k] - \text{expected})^2}{\text{expected}}
\end{align}

Large $\chi^2$ values indicate non-uniform distribution, suggesting correct key detection.
\end{examplebox}

\subsection{Implementation Attacks vs. Theoretical Attacks}

\begin{definitionbox}{Implementation Considerations}
Real-world differential attacks must consider:

\begin{enumerate}
    \item \textbf{Noise:} Random variations in measurements
    \item \textbf{Limited Data:} Finite number of plaintext-ciphertext pairs
    \item \textbf{Computational Complexity:} Time and memory constraints
    \item \textbf{False Positives:} Incorrect key candidates appearing correct
\end{enumerate}
\end{definitionbox}

\section{Case Studies: Historical Applications}

\subsection{DES Differential Cryptanalysis}

\begin{examplebox}{DES Attack Summary}
Biham and Shamir's attack on DES \cite{biham1993differential}:

\textbf{Characteristic:} 15-round differential with probability $2^{-47}$

\textbf{Data Complexity:} $2^{47}$ chosen plaintext pairs

\textbf{Time Complexity:} $2^{47}$ operations

\textbf{Success Rate:} Recovers 48 bits of the key

\textbf{Significance:} First demonstration that DES was weaker than brute force
\end{examplebox}

\subsection{Modern Cipher Analysis}

\begin{examplebox}{AES Differential Analysis}
AES differential analysis \cite{daemen2002design,nist2001aes} shows:

\textbf{4-round AES:} Best differential has probability $\approx 2^{-150}$

\textbf{Full AES-128:} No practical differential attack exists

\textbf{Design Success:} AES S-box and MixColumns provide excellent resistance

\textbf{Security Margin:} 10-14 rounds provide substantial protection
\end{examplebox}

\section{Tools and Software for Differential Analysis}

\subsection{Automated Characteristic Search}

Modern tools use sophisticated algorithms \cite{mouha2011automatic,sun2014automatic}:

\begin{enumerate}
    \item \textbf{SAT/SMT Solvers:} Model differential propagation as logical constraints
    \item \textbf{MILP (Mixed Integer Linear Programming):} Optimize characteristic probability
    \item \textbf{Machine Learning:} Neural networks for pattern recognition
    \item \textbf{Genetic Algorithms:} Evolutionary search for good characteristics
\end{enumerate}

\subsection{Educational Tools}

\begin{algorithmbox}{Simple Differential Analysis Tool}
\begin{algorithmic}[1]
\REQUIRE S-box definition, number of rounds, target cipher
\ENSURE Differential analysis results
\STATE Build differential distribution table for S-box
\STATE Generate all possible 1-round characteristics
\STATE Extend characteristics through multiple rounds
\STATE Calculate cumulative probabilities  
\STATE Rank characteristics by probability
\STATE Estimate attack complexity
\STATE Generate visual differential trails
\RETURN Analysis report with attack feasibility
\end{algorithmic}
\end{algorithmbox}

\section{Exercises and Practice Problems}

\subsection{Basic Exercises}

\begin{enumerate}
    \item \textbf{XOR Practice:} Compute differences for various bit strings
    \item \textbf{S-box Analysis:} Build DDT for a given 3-bit S-box
    \item \textbf{Probability Calculation:} Find differential probabilities through 2 rounds
    \item \textbf{Characteristic Extension:} Extend a 1-round characteristic to 3 rounds
\end{enumerate}

\subsection{Advanced Projects}

\begin{enumerate}
    \item \textbf{Mini-Cipher Implementation:} Implement and attack the educational cipher
    \item \textbf{S-box Optimization:} Design an S-box with minimal differential probability
    \item \textbf{Tool Development:} Create software for automatic characteristic search
    \item \textbf{Comparison Study:} Analyze multiple cipher designs for differential resistance
\end{enumerate}

\section{Conclusion and Further Reading}

\subsection{Key Takeaways}

Differential cryptanalysis provides:

\begin{enumerate}
    \item A systematic framework for analyzing symmetric ciphers
    \item Mathematical tools for quantifying cipher security
    \item Design principles for building resistant algorithms
    \item Foundation for advanced cryptanalytic techniques
\end{enumerate}

\subsection{Modern Developments}

Current research directions include:

\begin{enumerate}
    \item \textbf{Automated Tools:} AI-assisted cryptanalysis
    \item \textbf{Quantum Variants:} Quantum differential cryptanalysis
    \item \textbf{Post-Quantum Security:} Resistance in quantum computing era
    \item \textbf{Lightweight Cryptography:} Analysis of resource-constrained designs
\end{enumerate}

\subsection{Recommended Reading}

\begin{enumerate}
    \item Biham, E., \& Shamir, A. (1991). \textit{Differential Cryptanalysis of the Data Encryption Standard} \cite{biham1991differential}
    \item Matsui, M. (1993). \textit{Linear Cryptanalysis Method for DES Cipher} \cite{matsui1993linear}
    \item Daemen, J., \& Rijmen, V. (2002). \textit{The Design of Rijndael: AES} \cite{daemen2002design}
    \item Heys, H. M. (2002). \textit{A Tutorial on Linear and Differential Cryptanalysis} \cite{heys2002tutorial}
    \item Schneier, B. (1996). \textit{Applied Cryptography: Protocols, Algorithms, and Source Code in C} \cite{schneier1996applied}
    \item Menezes, A. J., Van Oorschot, P. C., \& Vanstone, S. A. (1996). \textit{Handbook of Applied Cryptography} \cite{menezes1996handbook}
    \item Stinson, D. R. (2005). \textit{Cryptography: Theory and Practice} \cite{stinson2005cryptography}
    \item Katz, J., \& Lindell, Y. (2014). \textit{Introduction to Modern Cryptography} \cite{katz2014introduction}
\end{enumerate}

\appendix

\section{Mathematical Reference}

\subsection{Probability Theory Essentials}

\begin{definition}[Conditional Probability]
For events $A$ and $B$ with $P(B) > 0$:
\begin{align}
P(A|B) = \frac{P(A \cap B)}{P(B)}
\end{align}
\end{definition}

\begin{definition}[Independence]
Events $A$ and $B$ are independent if:
\begin{align}
P(A \cap B) = P(A) \cdot P(B)
\end{align}
\end{definition}

\subsection{Information Theory}

\begin{definition}[Entropy]
For a random variable $X$ with probability mass function $p(x)$:
\begin{align}
H(X) = -\sum_{x} p(x) \log_2 p(x)
\end{align}
\end{definition}

\section{Implementation Code Examples}

\begin{lstlisting}[style=cryptoCode, caption={Differential Distribution Table (DDT) Construction and Analysis}]
#include <vector>
#include <array>

class DifferentialAnalyzer {
private:
    static constexpr size_t SBOX_SIZE = 16;
    std::array<uint8_t, SBOX_SIZE> sbox;
    std::array<std::array<int, SBOX_SIZE>, SBOX_SIZE> ddt;
    
public:
    void buildDDT() {
        // Initialize DDT to zeros
        for (auto& row : ddt) {
            row.fill(0);
        }
        
        // Fill DDT
        for (int dx = 0; dx < SBOX_SIZE; ++dx) {
            for (int x = 0; x < SBOX_SIZE; ++x) {
                int x_prime = x ^ dx;
                int dy = sbox[x] ^ sbox[x_prime];
                ddt[dx][dy]++;
            }
        }
    }
    
    double getDifferentialProbability(int input_diff, int output_diff) {
        return static_cast<double>(ddt[input_diff][output_diff]) / SBOX_SIZE;
    }
};
\end{lstlisting}

\begin{lstlisting}[style=cryptoCode, caption={Differential Characteristic Search and Validation}]
#include <iostream>
#include <iomanip>
#include <random>

class DifferentialCharacteristic {
private:
    std::vector<std::pair<uint16_t, uint16_t>> rounds;
    double totalProbability;
    
public:
    void addRound(uint16_t inputDiff, uint16_t outputDiff, 
                  double prob) {
        rounds.emplace_back(inputDiff, outputDiff);
        if (rounds.size() == 1) {
            totalProbability = prob;
        } else {
            totalProbability *= prob;
        }
    }
    
    // Validate characteristic with random pairs
    bool validateCharacteristic(const std::function<uint16_t(uint16_t)>& cipher,
                               int numPairs = 10000) {
        std::random_device rd;
        std::mt19937 gen(rd());
        std::uniform_int_distribution<uint16_t> dis;
        
        int matches = 0;
        uint16_t expectedInputDiff = rounds[0].first;
        uint16_t expectedOutputDiff = rounds.back().second;
        
        for (int i = 0; i < numPairs; ++i) {
            uint16_t plaintext1 = dis(gen);
            uint16_t plaintext2 = plaintext1 ^ expectedInputDiff;
            
            uint16_t ciphertext1 = cipher(plaintext1);
            uint16_t ciphertext2 = cipher(plaintext2);
            
            if ((ciphertext1 ^ ciphertext2) == expectedOutputDiff) {
                matches++;
            }
        }
        
        double observedProbability = 
            static_cast<double>(matches) / numPairs;
        double expectedProbability = totalProbability;
        
        // Check if observed probability is close to expected
        return std::abs(observedProbability - expectedProbability) 
               < 3 * std::sqrt(expectedProbability / numPairs);
    }
    
    void printCharacteristic() const {
        std::cout << "Differential Characteristic:\n";
        std::cout << "Total Probability: " << totalProbability << "\n";
        for (size_t i = 0; i < rounds.size(); ++i) {
            std::cout << "Round " << i + 1 << ": " 
                      << std::hex << rounds[i].first 
                      << " -> " << rounds[i].second << "\n";
        }
    }
};
\end{lstlisting}

\begin{lstlisting}[style=cryptoCode, caption={Key Recovery Attack Implementation}]
#include <map>
#include <algorithm>

class DifferentialAttack {
private:
    using PairData = std::pair<uint16_t, uint16_t>;
    std::vector<std::pair<PairData, PairData>> pairs;
    DifferentialCharacteristic characteristic;
    
public:
    void addPair(uint16_t p1, uint16_t c1, uint16_t p2, uint16_t c2) {
        pairs.emplace_back(std::make_pair(p1, c1), 
                          std::make_pair(p2, c2));
    }
    
    // Recover last round key using differential attack
    std::vector<uint16_t> recoverLastRoundKey(
        const std::function<uint16_t(uint16_t, uint16_t)>& lastRoundInverse,
        uint16_t expectedDiff) {
        
        std::map<uint16_t, int> keyVotes;
        
        // Try all possible subkey values
        for (uint32_t key = 0; key <= 0xFFFF; ++key) {
            int votes = 0;
            
            for (const auto& pair : pairs) {
                uint16_t c1 = pair.first.second;
                uint16_t c2 = pair.second.second;
                
                // Partially decrypt with candidate key
                uint16_t intermediate1 = lastRoundInverse(c1, key);
                uint16_t intermediate2 = lastRoundInverse(c2, key);
                
                // Check if difference matches expectation
                if ((intermediate1 ^ intermediate2) == expectedDiff) {
                    votes++;
                }
            }
            
            if (votes > 0) {
                keyVotes[key] = votes;
            }
        }
        
        // Sort candidates by votes
        std::vector<std::pair<uint16_t, int>> candidates(
            keyVotes.begin(), keyVotes.end());
        std::sort(candidates.begin(), candidates.end(),
                 [](const auto& a, const auto& b) {
                     return a.second > b.second;
                 });
        
        // Return top candidates
        std::vector<uint16_t> topKeys;
        for (size_t i = 0; i < std::min(candidates.size(), size_t(5)); ++i) {
            topKeys.push_back(candidates[i].first);
        }
        
        return topKeys;
    }
    
    void printStatistics() const {
        std::cout << "Attack Statistics:\n";
        std::cout << "Number of pairs: " << pairs.size() << "\n";
        std::cout << "Expected success rate: " 
                  << "Based on characteristic probability\n";
    }
};
\end{lstlisting}

\begin{thebibliography}{99}

\bibitem{biham1991differential}
Biham, E., \& Shamir, A. (1991).
\newblock \textit{Differential Cryptanalysis of the Data Encryption Standard}.
\newblock Springer-Verlag.

\bibitem{biham1993differential}
Biham, E., \& Shamir, A. (1993).
\newblock Differential cryptanalysis of DES-like cryptosystems.
\newblock \textit{Journal of Cryptology}, 4(1), 3-72.

\bibitem{matsui1993linear}
Matsui, M. (1993).
\newblock Linear cryptanalysis method for DES cipher.
\newblock In \textit{Advances in Cryptology - EUROCRYPT '93} (pp. 386-397). Springer.

\bibitem{daemen2002design}
Daemen, J., \& Rijmen, V. (2002).
\newblock \textit{The Design of Rijndael: AES - The Advanced Encryption Standard}.
\newblock Springer-Verlag.

\bibitem{heys2002tutorial}
Heys, H. M. (2002).
\newblock A tutorial on linear and differential cryptanalysis.
\newblock \textit{Cryptologia}, 26(3), 189-221.

\bibitem{lai1991markov}
Lai, X., Massey, J. L., \& Murphy, S. (1991).
\newblock Markov ciphers and differential cryptanalysis.
\newblock In \textit{Advances in Cryptology - EUROCRYPT '91} (pp. 17-38). Springer.

\bibitem{knudsen1994truncated}
Knudsen, L. R. (1994).
\newblock Truncated and higher order differentials.
\newblock In \textit{Fast Software Encryption} (pp. 196-211). Springer.

\bibitem{biham1999impossible}
Biham, E., Biryukov, A., \& Shamir, A. (1999).
\newblock Cryptanalysis of Skipjack reduced to 31 rounds using impossible differentials.
\newblock In \textit{Advances in Cryptology - EUROCRYPT '99} (pp. 12-23). Springer.

\bibitem{wagner1999boomerang}
Wagner, D. (1999).
\newblock The boomerang attack.
\newblock In \textit{Fast Software Encryption} (pp. 156-170). Springer.

\bibitem{kelsey1998amplified}
Kelsey, J., Kohno, T., \& Schneier, B. (1998).
\newblock Amplified boomerang attacks against reduced-round MARS and Serpent.
\newblock In \textit{Fast Software Encryption} (pp. 75-93). Springer.

\bibitem{mouha2011automatic}
Mouha, N., Wang, Q., Gu, D., \& Preneel, B. (2011).
\newblock Differential and linear cryptanalysis using mixed-integer linear programming.
\newblock In \textit{Information Security and Cryptology} (pp. 57-76). Springer.

\bibitem{sun2014automatic}
Sun, S., Hu, L., Wang, P., Qiao, K., Ma, X., \& Song, L. (2014).
\newblock Automatic security evaluation and (related-key) differential characteristic search: application to SIMON, PRESENT, LBlock, DES(L) and other bit-oriented block ciphers.
\newblock In \textit{Advances in Cryptology - ASIACRYPT 2014} (pp. 158-178). Springer.

\bibitem{kaliski1988linear}
Kaliski Jr, B. S., \& Robshaw, M. J. (1994).
\newblock Linear cryptanalysis using multiple approximations.
\newblock In \textit{Advances in Cryptology - CRYPTO '94} (pp. 26-39). Springer.

\bibitem{coppersmith1994data}
Coppersmith, D. (1994).
\newblock The Data Encryption Standard (DES) and its strength against attacks.
\newblock \textit{IBM Journal of Research and Development}, 38(3), 243-250.

\bibitem{nyberg1994linear}
Nyberg, K. (1994).
\newblock Linear approximation of block ciphers.
\newblock In \textit{Advances in Cryptology - EUROCRYPT '94} (pp. 439-444). Springer.

\bibitem{chabaud1998links}
Chabaud, F., \& Vaudenay, S. (1998).
\newblock Links between differential and linear cryptanalysis.
\newblock In \textit{Advances in Cryptology - EUROCRYPT '94} (pp. 356-365). Springer.

\bibitem{vaudenay1996decorrelation}
Vaudenay, S. (1996).
\newblock On the security of CS-cipher.
\newblock In \textit{Fast Software Encryption} (pp. 260-274). Springer.

\bibitem{aoki1998best}
Aoki, K., Ichikawa, T., Kanda, M., Matsui, M., Moriai, S., Nakajima, J., \& Tokita, T. (1998).
\newblock Camellia: A 128-bit block cipher suitable for multiple platforms.
\newblock In \textit{Selected Areas in Cryptography} (pp. 39-56). Springer.

\bibitem{schneier1996applied}
Schneier, B. (1996).
\newblock \textit{Applied Cryptography: Protocols, Algorithms, and Source Code in C}.
\newblock John Wiley \& Sons, 2nd edition.

\bibitem{menezes1996handbook}
Menezes, A. J., Van Oorschot, P. C., \& Vanstone, S. A. (1996).
\newblock \textit{Handbook of Applied Cryptography}.
\newblock CRC Press.

\bibitem{stinson2005cryptography}
Stinson, D. R. (2005).
\newblock \textit{Cryptography: Theory and Practice}.
\newblock Chapman and Hall/CRC, 3rd edition.

\bibitem{katz2014introduction}
Katz, J., \& Lindell, Y. (2014).
\newblock \textit{Introduction to Modern Cryptography}.
\newblock CRC Press, 2nd edition.

\bibitem{ferguson2010cryptography}
Ferguson, N., Schneier, B., \& Kohno, T. (2010).
\newblock \textit{Cryptography Engineering: Design Principles and Practical Applications}.
\newblock Wiley Publishing.

\bibitem{nist2001aes}
NIST (2001).
\newblock Advanced Encryption Standard (AES).
\newblock \textit{Federal Information Processing Standards Publication 197}.
\newblock U.S. Department of Commerce.

\bibitem{nist1999des}
NIST (1999).
\newblock Data Encryption Standard (DES).
\newblock \textit{Federal Information Processing Standards Publication 46-3}.
\newblock U.S. Department of Commerce.

\bibitem{iso2005block}
ISO/IEC 18033-3 (2005).
\newblock Information technology - Security techniques - Encryption algorithms - Part 3: Block ciphers.
\newblock International Organization for Standardization.

\bibitem{knudsen1998cryptanalysis}
Knudsen, L. R. (1998).
\newblock \textit{Contemporary Cryptology}.
\newblock In Contemporary Cryptology (pp. 209-270). Springer.

\bibitem{murphy1990probability}
Murphy, S. (1990).
\newblock The cryptanalysis of FEAL-4 with 20 chosen plaintexts.
\newblock \textit{Journal of Cryptology}, 2(3), 145-154.

\bibitem{gilbert1990practical}
Gilbert, H., \& Chassé, G. (1990).
\newblock A statistical attack of the FEAL-8 cryptosystem.
\newblock In \textit{Advances in Cryptology - CRYPTO '90} (pp. 22-33). Springer.

\bibitem{matsui1995security}
Matsui, M. (1995).
\newblock On correlation between the order of S-boxes and the strength of DES.
\newblock In \textit{Advances in Cryptology - EUROCRYPT '94} (pp. 366-375). Springer.

\bibitem{knudsen1995analysis}
Knudsen, L. R. (1995).
\newblock Cryptanalysis of LOKI91.
\newblock In \textit{Advances in Cryptology - AUSCRYPT '92} (pp. 196-208). Springer.

\end{thebibliography}

\end{document}
